\setcounter{section}{16}

  \zwischenueberschrift{Soal latihan}

\inputaufgabegibtloesung
{}
{

Kita meninjau parabola standar

\mavergleichskette
{\vergleichskette
{ Y
}
{ \subseteq }{ \R^2
}
{  }{
}
{    }{
}
{   }{
}
}
{}{}{}
dengan struktur Riemann terinduksi dan parametrisasi
\maabbeledisp {\varphi} {\R} { Y
} { t } { \left(  t   , \,  t^2                   \right)
} {,}
yang kita pandang sebagai suatu peta
\zusatzklammer {invers} {} {.}
Tentukan
\definitionsverweis {fungsi fundamental}{}{}
Riemann

\mavergleichskette
{\vergleichskette
{ g(t)
}
{ = }{ g_{11} (t)
}
{  }{
}
{    }{
}
{   }{
}
}
{}{}{.}

}
{} {}

\inputaufgabe
{}
{

Misalkan $M$ suatu
\definitionsverweis {manifold diferensiabel}{}{}
sedemikian sehingga
\definitionsverweis {bundel tangen}{}{}
$TM$ trivial. Tunjukkan bahwa melalui suatu trivialisasi diferensiabel

\mavergleichskette
{\vergleichskette
{ TM
}
{ \cong }{ M \times \R^n
}
{  }{
}
{    }{
}
{   }{
}
}
{}{}{}
\zusatzklammer {dengan bantuan
\definitionsverweis {hasil kali dalam standar}{}{}
pada $\R^n$} {} {}
$M$ dapat dijadikan suatu
\definitionsverweis {manifold Riemann}{}{.}

}
{} {}

\inputaufgabe
{}
{

Kita meninjau lingkaran

\mavergleichskette
{\vergleichskette
{ S^1
}
{ \subseteq }{ \R^2
}
{  }{
}
{    }{
}
{   }{
}
}
{}{}{}
dengan
\definitionsverweis {struktur Riemann terinduksi}{}{.}
Trivialisasi
\definitionsverweis {bundel tangen}{}{}
\maabbeledisp {} { S^1 \times \R } { T S^1
} {( a,b, u )} { (a,b, -ub,ua )
} {,}
melalui hasil kali dalam standar pada $\R$, juga menghasilkan suatu struktur Riemann pada $S^1$. Tunjukkan bahwa kedua struktur Riemann tersebut berimpit.

}
{} {}

\inputaufgabe
{}
{

Misalkan
\mathkor {} {L} {dan} {M} {}
merupakan
\definitionsverweis {manifold Riemann}{}{.}
Tunjukkan bahwa $L \times M$ secara alami juga merupakan suatu manifold Riemann.

}
{} {}

\inputaufgabe
{}
{

Kita meninjau suatu himpunan terbuka
\mathl{V \subseteq \R^n}{} sebagai
\definitionsverweis {manifold Riemann}{}{.} Apa
\definitionsverweis {bentuk volume kanonik}{}{} pada $V$?

}
{} {}

\inputaufgabe
{}
{

Misalkan $M$ suatu
\definitionsverweis {manifold Riemann}{}{.}
Tunjukkan bahwa
\definitionsverweis {pemetaan}{}{}
\maabbeledisp {} {
{ \mathcal V } ( M )  } {
{ \mathcal E }^{  1 } (  M   )
} { F } { \omega_F
} {,}
dengan

\mavergleichskettedisp
{\vergleichskette
{ \left(  \omega_F(P)  \right) (v)
}
{ \defeq} { \left\langle F(P) , v  \right\rangle_P
}
{ } {
}
{ } {
}
{ } {
}
}
{}{}{,}
bersifat
\definitionsverweis {linear}{}{.}

}
{} {}

\inputaufgabe
{}
{

Kita meninjau suatu himpunan terbuka
\mathl{V \subseteq \R^n}{} sebagai
\definitionsverweis {manifold Riemann}{}{.} Apa yang dinyatakan oleh korespondensi antara
\definitionsverweis {medan vektor}{}{}
dan
$1$-\definitionsverweis {bentuk diferensial}{}{}
yang diuraikan dalam
Lema 16.3
untuk situasi ini?

}
{} {}

\inputaufgabe
{}
{

Berikan alasan untuk
Catatan 16.4.

}
{} {}

\inputaufgabe
{}
{

Misalkan $M$ suatu
\definitionsverweis {manifold Riemann}{}{}
\definitionsverweis {terorientasi}{}{.} Tunjukkan bahwa
\definitionsverweis {bentuk volume kanonik}{}{} $\omega$ ditentukan oleh syarat bahwa pada setiap titik, nilainya pada setiap basis ortonormal yang merepresentasikan orientasi adalah $1$.

}
{} {}

\inputaufgabe
{}
{

Misalkan $V$ suatu ruang vektor riil terorientasi berdimensi $n$ dan $\lambda$ suatu
\definitionsverweis {ukuran yang invarian terhadap translasi}{}{}
pada $V$. Tunjukkan bahwa pemetaan
\maabbeledisp {} {V \times \cdots \times V} {\R
} {(v_1 , \ldots , v_n) } { \pm \lambda (P(v_1 , \ldots , v_n))
} {,}
dengan tanda positif dipilih jika vektor tersebut merepresentasikan orientasi, merupakan suatu pemetaan multilinear selang-seling.

}
{} {}

\inputaufgabe
{}
{

Tunjukkan bahwa pada suatu
\definitionsverweis {manifold Riemann}{}{},
\definitionsverweis {pemetaan koordinat}{}{}
\maabbdisp {\alpha} { U } { V
} {}
secara umum tidak menginduksi suatu
\definitionsverweis {isometri}{}{}
\maabbdisp {T_P(\alpha)} {T_PU} {T_{\alpha(P)} V
} {}
\zusatzklammer {jika
\mathl{T_PU}{} dilengkapi dengan
\mathl{\left\langle  - ,  -  \right\rangle_P}{} dan
\mathl{T_{\alpha(P)} V= \R^n}{} dilengkapi dengan hasil kali dalam standar} {} {.}

}
{} {}

\inputaufgabegibtloesung
{}
{

Kita meninjau grafik $M$ dari pemetaan
\maabbeledisp {\varphi} {\R^2} {\R
} {(u,v)} { u^2+uv-v^3
} {,} sebagai suatu submanifold tertutup berdimensi dua dari $\R^3$, yaitu

\mavergleichskettedisp
{\vergleichskette
{ M
}
{ =} { \left\{ (u,v,u^2+uv-v^3)  \mid  (u,v) \in \R^2         \right\}
}
{ } {
}
{ } {
}
{ } {
}
}
{}{}{,}
dengan metrik Riemann yang diinduksi dari $\R^3$. Pemetaan
\maabbeledisp {\psi} {\R^2} { M
} {(u,v)} { (u,v,u^2+uv-v^3)
} {,}
merupakan difeomorfisme yang terkait.

a) Tentukan diferensial total dari $\psi$ serta vektor citra
\mathl{T_P(\psi) (e_1)}{} dan
\mathl{T_P(\psi) (e_2)}{} dalam
\mathl{T_{\psi(P)}M}{.}

b) Untuk setiap titik berbentuk

\mavergleichskette
{\vergleichskette
{ P
}
{ = }{ (u,0)
}
{  }{
}
{    }{
}
{   }{
}
}
{}{}{}
tentukan luas jajaran genjang yang direntang oleh
\mathl{T_P(\psi) (e_1)}{} dan
\mathl{T_P(\psi) (e_2)}{} dalam
\mathl{T_{\psi(P)}M}{.}

c) Untuk setiap titik berbentuk

\mavergleichskette
{\vergleichskette
{ P
}
{ = }{ (0,v)
}
{  }{
}
{    }{
}
{   }{
}
}
{}{}{}
tentukan luas jajaran genjang yang direntang oleh
\mathl{T_P(\psi) (e_1)}{} dan
\mathl{T_P(\psi) (e_2)}{} dalam
\mathl{T_{\psi(P)}M}{.}

}
{} {}

\inputaufgabe
{}
{

Misalkan
\maabbdisp {\varphi} {\R^n } {\R
} {}
\zusatzklammer {dengan \mathlk{m=n -1 \geq 0}{}} {} {}
suatu
\definitionsverweis {fungsi yang terdiferensialkan secara kontinu}{}{,}
yang pada setiap titik dari
\definitionsverweis {serat}{}{}
$M$ di atas
\mathl{0 \in \R}{}, bersifat
\definitionsverweis {reguler}{}{.}
Kita memandang $M$ sebagai suatu manifold Riemann terorientasi. Tunjukkan bahwa antara bentuk volume $\tau$ dari
Korolari 15.6
dan
\definitionsverweis {bentuk volume kanonik}{}{}
$\omega$ berlaku hubungan

\mavergleichskettedisp
{\vergleichskette
{\tau(P,v_1 , \ldots , v_m)
}
{ =} { \pm \Vert {
\operatorname{Grad} \, \varphi ( P ) } \Vert  \omega(P, v_1 , \ldots , v_m)
}
{ } {
}
{ } {
}
{ } {
}
}
{}{}{.}

}
{} {}

\inputaufgabe
{}
{

Misalkan
\maabbdisp {\varphi} {\R^n } {\R^\ell
} {}
\zusatzklammer {dengan \mathlk{m=n - \ell \geq 0}{}} {} {}
suatu
\definitionsverweis {pemetaan yang terdiferensialkan secara kontinu}{}{,}
yang pada setiap titik dari
\definitionsverweis {serat}{}{}
$M$ di atas
\mathl{0 \in \R^\ell}{}, bersifat
\definitionsverweis {reguler}{}{.}
Kita memandang $M$ sebagai suatu manifold Riemann terorientasi. Andaikan bahwa gradien

\mathdisp {\operatorname{Grad} \, \varphi_1  (  P ) , \ldots ,
\operatorname{Grad} \, \varphi_\ell ( P )} {  }

saling tegak lurus pada setiap titik
\mathl{P\in M}{.} Tunjukkan bahwa antara bentuk volume $\tau$ dari
Korolari 15.6
dan
\definitionsverweis {bentuk volume kanonik}{}{}
$\omega$ berlaku hubungan

\mavergleichskettedisp
{\vergleichskette
{\tau(P,v_1 , \ldots , v_m)
}
{ =} { \pm \Vert {
\operatorname{Grad} \, \varphi_1  (  P ) } \Vert  \cdots \Vert {
\operatorname{Grad} \, \varphi_\ell ( P ) } \Vert \omega(P, v_1 , \ldots , v_m)
}
{ } {
}
{ } {
}
{ } {
}
}
{}{}{.}

}
{} {}

\inputaufgabe
{}
{

Tunjukkan bahwa
\definitionsverweis {bentuk luas kanonik}{}{}
pada
\definitionsverweis {permukaan bola satuan}{}{}

\mavergleichskette
{\vergleichskette
{S^2
}
{ \subset }{ \R^3
}
{  }{
}
{    }{
}
{   }{
}
}
{}{}{}
adalah

\mathdisp {x dy \wedge dz  -y dx \wedge dz +  z dx \wedge dy} {  }

dengan $x,y,z$ menyatakan koordinat pada $\R^3$.

}
{} {}

Suatu
\definitionsverweis {bundel vektor}{}{}
\maabb {p} { E } { X
} {}
di atas suatu
\definitionsverweis {manifold diferensiabel}{}{}
$X$ disebut
\definitionswort {bundel vektor Riemann}{,}
jika pada setiap serat $E_x$ diberikan suatu
\definitionsverweis {hasil kali dalam}{}{}
dengan sifat bahwa terdapat
\definitionsverweis {peta}{}{}
yang mentrivialkan
\maabbdisp {\alpha} { U } { V \subseteq \R^n
} {}
bersama
\maabbdisp {\theta} {E\vert_U \cong U \times \R^r } { V \times \R^r
} {}
sedemikian sehingga fungsi

\mavergleichskettedisp
{\vergleichskette
{ h_{ij} (Q )
}
{ \defeq} { \left\langle  \theta^{-1}(Q,e_i)  , \theta^{-1}(Q,e_j)    \right\rangle_{\alpha^{-1}(Q)}
}
{ } {
}
{ } {
}
{ } {
}
}
{}{}{}
yang terdiferensialkan secara kontinu.

Jadi, pada suatu
\definitionsverweis {manifold Riemann}{}{},
\definitionsverweis {bundel tangen}{}{}
merupakan suatu bundel vektor Riemann.

\inputaufgabe
{}
{

Misalkan
\maabb {p} { E } { M
} {}
suatu
\definitionsverweis {bundel vektor Riemann}{}{}
di atas suatu
\definitionsverweis {manifold diferensiabel}{}{}
$M$, dan misalkan
\maabb {\varphi} { L } { M
} {}
suatu
\definitionsverweis {pemetaan diferensiabel}{}{}
dengan
\definitionsverweis {bundel vektor tarik balik}{}{}
\maabb {} {\varphi^*E} {L
} {.}
Tunjukkan bahwa $\varphi^*E$ secara alami juga merupakan suatu bundel vektor Riemann.

}
{} {}

  \zwischenueberschrift{Soal untuk dikumpulkan}

Pada suatu manifold Riemann $M$, untuk suatu vektor singgung

\mavergleichskette
{\vergleichskette
{ v
}
{ \in }{ T_PM
}
{  }{
}
{    }{
}
{   }{
}
}
{}{}{}
normanya didefinisikan oleh

\mavergleichskette
{\vergleichskette
{ \Vert {v} \Vert
}
{ = }{ \sqrt{ \left\langle v , v  \right\rangle_P }
}
{  }{
}
{    }{
}
{   }{
}
}
{}{}{.}

\inputaufgabe
{4}
{

Misalkan $M$ suatu
\definitionsverweis {manifold Riemann}{}{.}
Tunjukkan bahwa pemetaan
\maabbeledisp {} {TM} {\R
} { v } { \Vert { v } \Vert
} {,}
bersifat
\definitionsverweis {kontinu}{}{.}

}
{} {}

\inputaufgabe
{3}
{

Tunjukkan bahwa
\mathl{\R \times \R_+}{} dengan bentuk bilinear yang diberikan oleh
\definitionsverweis {bentuk Hesse}{}{} dari
\definitionsverweis {fungsi}{}{}
\maabbeledisp {f} {\R \times \R_+} {\R
} {(x,y)} { x^2+y^4
} {,}
merupakan suatu
\definitionsverweis {manifold Riemann}{}{.}

}
{} {}

\inputaufgabe
{4}
{

Untuk setiap titik
\mathl{P=(x,y,z)}{} pada
\definitionsverweis {permukaan bola satuan}{}{}
$K$, berikan suatu
\definitionsverweis {basis ortonormal}{}{}
dalam
\mathl{T_PK \subset \R^3}{}
\zusatzklammer {terhadap
\definitionsverweis {struktur Riemann terinduksi}{}{}} {} {.}

}
{} {}

\inputaufgabe
{6}
{

Dalam $\R^3$, diberikan
\definitionsverweis {elipsoid}{}{}

\mathdisp {E= \left\{ (x,y,z)  \mid   x^2+y^2+3z^2 \leq 5        \right\}} {  }
dan bidang

\mathdisp {M= \left\{ (x,y,z)  \mid   7x-3y-2z =  2         \right\}} {  }.
Hitunglah luas irisan
\mathl{M \cap E}{.}

}
{} {}

\inputaufgabe
{6}
{

Buatlah suatu grafik komputer yang memvisualisasikan situasi yang diuraikan dalam
Catatan 16.4
dengan menggunakan suatu permukaan di $\R^3$.

}
{} {}
